We verify every primitive defining the class. The full-data vectors are i.i.d. because the two Rademacher potential outcomes are independent across occasions and arms and the context is constant. The displayed policy is the logged assignment rule under \(P^{\dagger}\); at each occasion it is a function of the past-measurable state, and the action randomization is independent of the current potential outcomes. Thus the i.i.d.-arrival and sequential-randomization assumptions hold.

At the two regime representatives the arm-probability vectors are
\[
 (2/5,1/4,7/20),
 \qquad (2/5,1/6,13/30). \tag{1}
\]
Linear interpolation preserves their coordinatewise minima, so every arm probability is at least \(1/6>\varepsilon\). The coordinate slopes are \(0,-1/12,1/12\), proving the policy-Lipschitz condition with \(L_\Pi=1/12\).

For \(\theta\in[-1,1]^2\),
\[
 \sup|Y_t(1)-\theta_1|\le1+\sqrt{2/5}<2,
 \qquad
 \sup|Y_t(2)-\theta_2|\le1+\sqrt{1/2}<2. \tag{2}
\]
Each score has only one possibly nonzero coordinate, so \(B=2\) is a valid common envelope. The derivative matrices are \(\operatorname{diag}(-1,0)\), \(\operatorname{diag}(0,-1)\), and zero. Their operator norms are at most one, and every second derivative vanishes. Thus \(D_1=D_2=1\) are valid class envelopes, with actual curvature zero. Direct expectation gives
\[
 M_{P^{\dagger}}(\theta)=(-\theta_1,-\theta_2)^\top=-\theta,
 \qquad H=-I_2. \tag{3}
\]
It follows that \(\theta^*=0\) is the unique root, \(\|M_{P^{\dagger}}(\theta)\|_2=\|\theta\|_2\), \(\kappa(u)=u\) is valid, \(\operatorname{dist}(0,\partial[-1,1]^2)=1\), and \(\sigma_{\min}(H)=1\).

Every realized empirical moment has a zero in \(\Theta\). If arm \(a\in\{1,2\}\) is observed, choose \(\theta_a\) to be the average of its observed outcomes with positive weights \(1/\pi_t(a)\). This value lies in the convex hull of the two possible outcomes for that arm and hence in \([-1,1]\). If the arm is never observed, its empirical coordinate moment is identically zero and any coordinate in \([-1,1]\) works. The two coordinates separate, so these choices make \(G_T(\theta)=0\). Therefore the empirical moment norm has exact minimum zero.

At the first occasion the state is \(1/2\), and the arm-one probability is \(2/5\). Hence
\[
 P^{\dagger}(J=1)=P^{\dagger}(A_1=1)=2/5,
 \qquad P^{\dagger}(J=2)=3/5. \tag{4}
\]
The regime label is thus determined after the first assignment and has positive-mass support \(\mathcal J=\{1,2\}\). For every \(t\ge1\), \(\eta_t=\bar\eta_J\), while \(|\eta_0-\bar\eta_J|=1/2\). Consequently
\[
 T^{-1}\sum_{t=1}^T|\eta_{t-1}-\bar\eta_J|=\frac1{2T} \tag{5}
\]
almost surely. The terminal-measurability condition holds because \(J\) is \(\mathcal F_1\)-measurable; the uniform Cesaro condition holds with \(r_T=1/(2T)\) and \(\delta_T=0\); and \(e_T=T^{-1/2}+1/(2T)\to0\). Moreover,
\[
 0<\rho=\varepsilon/8<\min\{1,\varepsilon/(4D_2)\}, \tag{6}
\]
so the prescribed root-search radius is legal.

At \(\theta^*=0\), only arms one and two contribute to the covariance. Since \(E[Y_t(1)^2]=2/5\) and \(E[Y_t(2)^2]=1/2\),
\[
 \Omega_1
 =\operatorname{diag}\left(\frac{2/5}{2/5},\frac{1/2}{1/4}\right)
 =\operatorname{diag}(1,2), \tag{7}
\]
and
\[
 \Omega_2
 =\operatorname{diag}\left(\frac{2/5}{2/5},\frac{1/2}{1/6}\right)
 =\operatorname{diag}(1,3). \tag{8}
\]
Both matrices dominate \(I_2\), so the regime-nondegeneracy assumption holds with \(\omega_-=1\). Because \(H=-I_2\), the sandwich definition gives \(V_j=\Omega_j\). Equations (1)--(8) verify every class condition with constants independent of \(T\).

Finally, for \(\Sigma=\operatorname{diag}(1,2)\) and \(c=(1,0)^\top\),
\[
 c^\top V_1c=c^\top V_2c=c^\top\Sigma c=1,
 \qquad V_1\ne V_2. \tag{9}
\]
The constant sequence has fixed weights \((2/5,3/5)\) and fixed covariance matrices. The contrastwise theorem applied to (9) therefore gives every-level validity for this contrast, while the matrixwise theorem rules out validity of the same deterministic matrix for every contrast. This proves the claimed in-class separation without making any assertion about dynamical basins of the Guo--Xu recursion.